\chapter{Polynômes, racines et factorisation dans le plan complexe}
\section*{Exercice 27.1 -- Second degré réel}
[Base] Résoudre $z^2-2z+5=0$ dans $\mathbb C$.

\section*{Exercice 27.2 -- Factorisation}
[Méthode] Factoriser $P(z)=z^3-1$ puis résoudre $P(z)=0$.

\section*{Exercice 27.3 -- Racine double}
[Raisonnement] Factoriser $z^3-3z+2$ sachant que 1 est racine. Combien y a-t-il de racines distinctes ?

\section*{Exercice 27.4 -- Racines conjuguées}
[Preuve] Montrer que si un polynôme à coefficients réels a pour racine $a+i b$, alors $a-i b$ est aussi racine.

\section*{Exercice 27.5 -- Nombre de racines}
[Preuve] Pourquoi un polynôme de degré au plus 4 s'annulant en cinq points distincts est-il identiquement nul ?

\section*{Exercice 27.6 -- Racine complexe d’un complexe}
[Approfondissement] Résoudre $z^2=3+4i$ en posant $z=x+i y$.
